% --- Archon physics formalization source begin ---
% archon:physics
% archon:covers PhyXMiniProblems/problem_phyx_mini_0732.lean
% archon:source-report reports/phyx_mini/problem_phyx_mini_0732.source.json

\chapter{Physics problem phyx\_mini\_0732}
\label{ch:PhyXMiniProblems_problem_phyx_mini_0732}

\paragraph{Problem source.}
\#\# Physical scenario

Figure shows part of a street where traffic flow is to be controlled to allow a platoon of cars to move smoothly along the street. Suppose that the platoon leaders have just reached intersection 2, where the green appeared when they were distance \$d\$ from the intersection. They continue to travel at a certain speed \$v\_p\$ (the speed limit) to reach intersection 3, where the green appears when they are distance \$d\$ from it. The intersections are separated by distances \$D\_\{23\}\$ and \$D\_\{12\}\$.

\#\# Question

What should be the time delay of the onset of green at intersection 3 relative to that at intersection 2 to keep the platoon moving smoothly?

\#\# Answer choices

- A: A: \$D\_\{13\} / v\_p\$

- B: B: \$2 D\_\{23\} / v\_p\$

- C: C: \$D\_\{23\} / v\_p\$

- D: D: \$D\_\{23\} / 2 v\_p\$

\#\# Figure caption (auxiliary; use the image as primary evidence)

The image depicts a traffic scene with a one-way street, indicated by a black sign with a white arrow pointing to the right, labeled "ONE WAY."

- **Background:**
  - The street runs horizontally from left to right.
  - There are three vertical intersecting roads, labeled with the numbers 1, 2, and 3 in sequence from left to right.

- **Vehicles:**
  - On the left side, there is a series of cars in different colors: black, white, yellow, red, and green. They are driving on the left lane of a two-lane road.

- **Measurements:**
  - Distances between intersections are marked with arrows labeled \textbackslash{}(D\_\{12\}\textbackslash{}) and \textbackslash{}(D\_\{23\}\textbackslash{}), denoting the distances between intersections 1 and 2, and intersections 2 and 3, respectively.

The intersections, vehicles, and directional sign convey a structured traffic layout typical for urban environments.

\#\# Dataset metadata

Kinematics; Multi-Formula Reasoning, Predictive Reasoning

\paragraph{Recorded answer/context.}
C: C: \$D\_\{23\} / v\_p\$

\paragraph{Figure/image path.}
/root/proposal\_for\_physic/science-mango/phyx\_mini\_run/phyx\_data/test\_image/732.png

\paragraph{Formalization target.}
create a compiling Lean file with sorry bodies at `PhyXMiniProblems/problem\_phyx\_mini\_0732.lean`. The Lean declarations must preserve the physical quantities, dimensions or dimensional roles, figure labels, governing-law hypotheses, and final relation expressed by this problem.
Use Mathlib/Physlib names found through LeanExplore where available. If a physics API is missing, introduce faithful local abstractions rather than scalar placeholder aliases.

\begin{theorem}[Physics formalization target]
\label{thm:physics:phyx_mini_0732:target}
\lean{PhyXMiniProblems.ProblemPhyXMini0732.problem_phyx_mini_0732}
\uses{def:physics:phyx-mini-0732:phyxminiproblems-problemphyxmini0732-lengthreadout, def:physics:phyx-mini-0732:phyxminiproblems-problemphyxmini0732-speedreadout, def:physics:phyx-mini-0732:phyxminiproblems-problemphyxmini0732-trafficsignalsetup, def:physics:phyx-mini-0732:phyxminiproblems-problemphyxmini0732-matchesprimarytrafficfigure, def:physics:phyx-mini-0732:phyxminiproblems-problemphyxmini0732-matchesplatoonscenario, def:physics:phyx-mini-0732:phyxminiproblems-problemphyxmini0732-satisfiesstreetdistancegeometry, def:physics:phyx-mini-0732:phyxminiproblems-problemphyxmini0732-matchesgreentriggerprotocol, def:physics:phyx-mini-0732:phyxminiproblems-problemphyxmini0732-hasphysicaltrafficparameters, def:physics:phyx-mini-0732:phyxminiproblems-problemphyxmini0732-satisfiesuniformplatoonmotion, def:physics:phyx-mini-0732:phyxminiproblems-problemphyxmini0732-greenonsetdelayreadout}
The assigned autoformalize agent should translate this physics problem into Lean declarations in the covered file, with theorem and lemma proof bodies written as `by sorry`.
\end{theorem}
\begin{proof}
This is an autoformalization task, not a proof task. Produce statements that can later be proved by the physics prover without weakening the physical model.
\end{proof}

% --- Archon declaration topology begin ---
\section{Declaration topology}
\begin{definition}[Length Quantity]
  \label{def:physics:phyx-mini-0732:phyxminiproblems-problemphyxmini0732-lengthquantity}
  \lean{PhyXMiniProblems.ProblemPhyXMini0732.LengthQuantity}
  A unit-independent physical length or signed coordinate along the street
\end{definition}
\begin{proof}
  This is the declared model definition of Length Quantity.
\end{proof}
\begin{definition}[Time Quantity]
  \label{def:physics:phyx-mini-0732:phyxminiproblems-problemphyxmini0732-timequantity}
  \lean{PhyXMiniProblems.ProblemPhyXMini0732.TimeQuantity}
  A unit-independent physical time coordinate
\end{definition}
\begin{proof}
  This is the declared model definition of Time Quantity.
\end{proof}
\begin{definition}[Speed Quantity]
  \label{def:physics:phyx-mini-0732:phyxminiproblems-problemphyxmini0732-speedquantity}
  \lean{PhyXMiniProblems.ProblemPhyXMini0732.SpeedQuantity}
  The nonnegative, unit-independent physical speed type supplied by Physlib
\end{definition}
\begin{proof}
  This is the declared model definition of Speed Quantity.
\end{proof}
\begin{definition}[length Readout]
  \label{def:physics:phyx-mini-0732:phyxminiproblems-problemphyxmini0732-lengthreadout}
  \lean{PhyXMiniProblems.ProblemPhyXMini0732.lengthReadout}
  \uses{def:physics:phyx-mini-0732:phyxminiproblems-problemphyxmini0732-lengthquantity}
  Read a length or longitudinal coordinate in the selected length unit
\end{definition}
\begin{proof}
  This is the declared model definition of length Readout.
\end{proof}
\begin{definition}[time Readout]
  \label{def:physics:phyx-mini-0732:phyxminiproblems-problemphyxmini0732-timereadout}
  \lean{PhyXMiniProblems.ProblemPhyXMini0732.timeReadout}
  \uses{def:physics:phyx-mini-0732:phyxminiproblems-problemphyxmini0732-timequantity}
  Read a time coordinate in the selected time unit
\end{definition}
\begin{proof}
  This is the declared model definition of time Readout.
\end{proof}
\begin{definition}[speed Readout]
  \label{def:physics:phyx-mini-0732:phyxminiproblems-problemphyxmini0732-speedreadout}
  \lean{PhyXMiniProblems.ProblemPhyXMini0732.speedReadout}
  \uses{def:physics:phyx-mini-0732:phyxminiproblems-problemphyxmini0732-speedquantity}
  Read a speed in the coherent selected length and time units
\end{definition}
\begin{proof}
  This is the declared model definition of speed Readout.
\end{proof}
\begin{definition}[Intersection]
  \label{def:physics:phyx-mini-0732:phyxminiproblems-problemphyxmini0732-intersection}
  \lean{PhyXMiniProblems.ProblemPhyXMini0732.Intersection}
  The three numbered intersections in their left-to-right order
\end{definition}
\begin{proof}
  This is the declared model definition of Intersection.
\end{proof}
\begin{definition}[Street Segment]
  \label{def:physics:phyx-mini-0732:phyxminiproblems-problemphyxmini0732-streetsegment}
  \lean{PhyXMiniProblems.ProblemPhyXMini0732.StreetSegment}
  The two consecutive street segments carrying the labels D₁₂ and D₂₃
\end{definition}
\begin{proof}
  This is the declared model definition of Street Segment.
\end{proof}
\begin{definition}[start Intersection]
  \label{def:physics:phyx-mini-0732:phyxminiproblems-problemphyxmini0732-streetsegment-startintersection}
  \lean{PhyXMiniProblems.ProblemPhyXMini0732.StreetSegment.startIntersection}
  \uses{def:physics:phyx-mini-0732:phyxminiproblems-problemphyxmini0732-intersection, def:physics:phyx-mini-0732:phyxminiproblems-problemphyxmini0732-streetsegment}
  The left endpoint of a directed street segment
\end{definition}
\begin{proof}
  This is the declared model definition of start Intersection.
\end{proof}
\begin{definition}[finish Intersection]
  \label{def:physics:phyx-mini-0732:phyxminiproblems-problemphyxmini0732-streetsegment-finishintersection}
  \lean{PhyXMiniProblems.ProblemPhyXMini0732.StreetSegment.finishIntersection}
  \uses{def:physics:phyx-mini-0732:phyxminiproblems-problemphyxmini0732-intersection, def:physics:phyx-mini-0732:phyxminiproblems-problemphyxmini0732-streetsegment}
  The right endpoint of a directed street segment
\end{definition}
\begin{proof}
  This is the declared model definition of finish Intersection.
\end{proof}
\begin{definition}[Segment Distance Label]
  \label{def:physics:phyx-mini-0732:phyxminiproblems-problemphyxmini0732-segmentdistancelabel}
  \lean{PhyXMiniProblems.ProblemPhyXMini0732.SegmentDistanceLabel}
  Literal distance labels printed below the street in the primary figure
\end{definition}
\begin{proof}
  This is the declared model definition of Segment Distance Label.
\end{proof}
\begin{definition}[expected Distance Label]
  \label{def:physics:phyx-mini-0732:phyxminiproblems-problemphyxmini0732-streetsegment-expecteddistancelabel}
  \lean{PhyXMiniProblems.ProblemPhyXMini0732.StreetSegment.expectedDistanceLabel}
  \uses{def:physics:phyx-mini-0732:phyxminiproblems-problemphyxmini0732-streetsegment, def:physics:phyx-mini-0732:phyxminiproblems-problemphyxmini0732-segmentdistancelabel}
  The printed label belonging to each directed segment
\end{definition}
\begin{proof}
  This is the declared model definition of expected Distance Label.
\end{proof}
\begin{definition}[Horizontal Direction]
  \label{def:physics:phyx-mini-0732:phyxminiproblems-problemphyxmini0732-horizontaldirection}
  \lean{PhyXMiniProblems.ProblemPhyXMini0732.HorizontalDirection}
  Direction of travel along the horizontal street axis
\end{definition}
\begin{proof}
  This is the declared model definition of Horizontal Direction.
\end{proof}
\begin{definition}[Answer Choice]
  \label{def:physics:phyx-mini-0732:phyxminiproblems-problemphyxmini0732-answerchoice}
  \lean{PhyXMiniProblems.ProblemPhyXMini0732.AnswerChoice}
  Labels of the four displayed answers
\end{definition}
\begin{proof}
  This is the declared model definition of Answer Choice.
\end{proof}
\begin{definition}[Traffic Control Figure]
  \label{def:physics:phyx-mini-0732:phyxminiproblems-problemphyxmini0732-trafficcontrolfigure}
  \lean{PhyXMiniProblems.ProblemPhyXMini0732.TrafficControlFigure}
  \uses{def:physics:phyx-mini-0732:phyxminiproblems-problemphyxmini0732-intersection, def:physics:phyx-mini-0732:phyxminiproblems-problemphyxmini0732-streetsegment, def:physics:phyx-mini-0732:phyxminiproblems-problemphyxmini0732-segmentdistancelabel, def:physics:phyx-mini-0732:phyxminiproblems-problemphyxmini0732-horizontaldirection}
  A structured transcription of the primary bitmap. It stores only discrete labels and qualitative visible features; physical distances are kept in the setup below
\end{definition}
\begin{proof}
  This is the declared model definition of Traffic Control Figure.
\end{proof}
\begin{definition}[Traffic Signal Setup]
  \label{def:physics:phyx-mini-0732:phyxminiproblems-problemphyxmini0732-trafficsignalsetup}
  \lean{PhyXMiniProblems.ProblemPhyXMini0732.TrafficSignalSetup}
  \uses{def:physics:phyx-mini-0732:phyxminiproblems-problemphyxmini0732-lengthquantity, def:physics:phyx-mini-0732:phyxminiproblems-problemphyxmini0732-timequantity, def:physics:phyx-mini-0732:phyxminiproblems-problemphyxmini0732-speedquantity, def:physics:phyx-mini-0732:phyxminiproblems-problemphyxmini0732-intersection, def:physics:phyx-mini-0732:phyxminiproblems-problemphyxmini0732-streetsegment, def:physics:phyx-mini-0732:phyxminiproblems-problemphyxmini0732-horizontaldirection, def:physics:phyx-mini-0732:phyxminiproblems-problemphyxmini0732-trafficcontrolfigure}
  Independent physical quantities and event observables. The green-onset times and leader trajectory are free fields here; their relations are imposed by separate protocol and kinematics predicates
\end{definition}
\begin{proof}
  This is the declared model definition of Traffic Signal Setup.
\end{proof}
\begin{definition}[Matches Primary Traffic Figure]
  \label{def:physics:phyx-mini-0732:phyxminiproblems-problemphyxmini0732-matchesprimarytrafficfigure}
  \lean{PhyXMiniProblems.ProblemPhyXMini0732.MatchesPrimaryTrafficFigure}
  \uses{def:physics:phyx-mini-0732:phyxminiproblems-problemphyxmini0732-trafficsignalsetup}
  Qualitative labels and geometry read directly from 732.png
\end{definition}
\begin{proof}
  This is the declared model definition of Matches Primary Traffic Figure.
\end{proof}
\begin{definition}[Matches Platoon Scenario]
  \label{def:physics:phyx-mini-0732:phyxminiproblems-problemphyxmini0732-matchesplatoonscenario}
  \lean{PhyXMiniProblems.ProblemPhyXMini0732.MatchesPlatoonScenario}
  \uses{def:physics:phyx-mini-0732:phyxminiproblems-problemphyxmini0732-timereadout, def:physics:phyx-mini-0732:phyxminiproblems-problemphyxmini0732-trafficsignalsetup}
  The prose says that the leaders have reached intersection 2, continue in the positive street direction, and travel at the speed limit. The earlier green onset at intersection 2 remains a separate event
\end{definition}
\begin{proof}
  This is the declared model definition of Matches Platoon Scenario.
\end{proof}
\begin{definition}[Satisfies Street Distance Geometry]
  \label{def:physics:phyx-mini-0732:phyxminiproblems-problemphyxmini0732-satisfiesstreetdistancegeometry}
  \lean{PhyXMiniProblems.ProblemPhyXMini0732.SatisfiesStreetDistanceGeometry}
  \uses{def:physics:phyx-mini-0732:phyxminiproblems-problemphyxmini0732-lengthreadout, def:physics:phyx-mini-0732:phyxminiproblems-problemphyxmini0732-streetsegment, def:physics:phyx-mini-0732:phyxminiproblems-problemphyxmini0732-trafficsignalsetup}
  The marked physical segment length is the difference of its endpoint coordinates. This records both D₁₂ and D₂₃, even though only D₂₃ survives in the requested delay
\end{definition}
\begin{proof}
  This is the declared model definition of Satisfies Street Distance Geometry.
\end{proof}
\begin{definition}[Matches Green Trigger Protocol]
  \label{def:physics:phyx-mini-0732:phyxminiproblems-problemphyxmini0732-matchesgreentriggerprotocol}
  \lean{PhyXMiniProblems.ProblemPhyXMini0732.MatchesGreenTriggerProtocol}
  \uses{def:physics:phyx-mini-0732:phyxminiproblems-problemphyxmini0732-lengthreadout, def:physics:phyx-mini-0732:phyxminiproblems-problemphyxmini0732-intersection, def:physics:phyx-mini-0732:phyxminiproblems-problemphyxmini0732-trafficsignalsetup}
  At each relevant green onset, the leader is the same physical distance d upstream of that intersection. These are trigger-event readouts, not a delay formula
\end{definition}
\begin{proof}
  This is the declared model definition of Matches Green Trigger Protocol.
\end{proof}
\begin{definition}[Has Physical Traffic Parameters]
  \label{def:physics:phyx-mini-0732:phyxminiproblems-problemphyxmini0732-hasphysicaltrafficparameters}
  \lean{PhyXMiniProblems.ProblemPhyXMini0732.HasPhysicalTrafficParameters}
  \uses{def:physics:phyx-mini-0732:phyxminiproblems-problemphyxmini0732-lengthreadout, def:physics:phyx-mini-0732:phyxminiproblems-problemphyxmini0732-timereadout, def:physics:phyx-mini-0732:phyxminiproblems-problemphyxmini0732-speedreadout, def:physics:phyx-mini-0732:phyxminiproblems-problemphyxmini0732-trafficsignalsetup}
  Nondegeneracy and event ordering for the physical branch depicted in the problem. No magnitude for the requested delay occurs here
\end{definition}
\begin{proof}
  This is the declared model definition of Has Physical Traffic Parameters.
\end{proof}
\begin{definition}[Satisfies Uniform Platoon Motion]
  \label{def:physics:phyx-mini-0732:phyxminiproblems-problemphyxmini0732-satisfiesuniformplatoonmotion}
  \lean{PhyXMiniProblems.ProblemPhyXMini0732.SatisfiesUniformPlatoonMotion}
  \uses{def:physics:phyx-mini-0732:phyxminiproblems-problemphyxmini0732-timequantity, def:physics:phyx-mini-0732:phyxminiproblems-problemphyxmini0732-lengthreadout, def:physics:phyx-mini-0732:phyxminiproblems-problemphyxmini0732-timereadout, def:physics:phyx-mini-0732:phyxminiproblems-problemphyxmini0732-speedreadout, def:physics:phyx-mini-0732:phyxminiproblems-problemphyxmini0732-trafficsignalsetup}
  Governing uniform-motion law during the interval from the onset of green at intersection 2 through the onset at intersection 3. It is deliberately quantified over arbitrary event times in that interval, rather than assuming the requested endpoint-delay formula
\end{definition}
\begin{proof}
  This is the declared model definition of Satisfies Uniform Platoon Motion.
\end{proof}
\begin{definition}[green Onset Delay Readout]
  \label{def:physics:phyx-mini-0732:phyxminiproblems-problemphyxmini0732-greenonsetdelayreadout}
  \lean{PhyXMiniProblems.ProblemPhyXMini0732.greenOnsetDelayReadout}
  \uses{def:physics:phyx-mini-0732:phyxminiproblems-problemphyxmini0732-timereadout, def:physics:phyx-mini-0732:phyxminiproblems-problemphyxmini0732-trafficsignalsetup}
  Scalar readout of the onset delay of green 3 relative to green 2
\end{definition}
\begin{proof}
  This is the declared model definition of green Onset Delay Readout.
\end{proof}
\begin{definition}[displayed Delay Readout]
  \label{def:physics:phyx-mini-0732:phyxminiproblems-problemphyxmini0732-displayeddelayreadout}
  \lean{PhyXMiniProblems.ProblemPhyXMini0732.displayedDelayReadout}
  \uses{def:physics:phyx-mini-0732:phyxminiproblems-problemphyxmini0732-lengthreadout, def:physics:phyx-mini-0732:phyxminiproblems-problemphyxmini0732-speedreadout, def:physics:phyx-mini-0732:phyxminiproblems-problemphyxmini0732-answerchoice, def:physics:phyx-mini-0732:phyxminiproblems-problemphyxmini0732-trafficsignalsetup}
  The delay expression printed beside an answer label. Choice A uses D₁₃ = D₁₂ + D₂₃, while choices B, C, and D use the stated multiples of D₂₃
\end{definition}
\begin{proof}
  This is the declared model definition of displayed Delay Readout.
\end{proof}
% --- Archon declaration topology end ---
% --- Archon physics formalization source end ---
